% Created 2020-10-25 So 23:03
% Intended LaTeX compiler: pdflatex
\UseRawInputEncoding

\documentclass[11pt]{scrartcl}
\usepackage[utf8]{inputenc}
\usepackage[T1]{fontenc}
\usepackage[ngerman]{babel}
\usepackage{graphicx}
\usepackage{grffile}
\usepackage{longtable}
\usepackage{wrapfig}
\usepackage{rotating}
\usepackage[normalem]{ulem}
\usepackage{amsmath}
\usepackage{trfsigns}
\usepackage{textcomp}
\usepackage{amssymb}
\usepackage{capt-of}
\usepackage{MnSymbol}
\usepackage{mathtools}
\usepackage{setspace}
\usepackage{nicematrix}
\usepackage{pdfpages}
% \usepackage{getlab}
\usepackage{listings}
\usepackage{pdfpages}
\usepackage[straightvoltages, european resistors, american inductors]{circuitikz}
\usetikzlibrary{arrows.meta}
\tikzset{FARROW/.style={arrows={-Triangle[angle=45:2.0mm]}}}
%\usepackage{hyperref}
% \usepackage{varwidth}


\definecolor{darkspringgreen}{rgb}{0.09, 0.45, 0.27}    % Farbe für die Kommentare bei Listings
\lstset{
  language= Matlab,                     % Setzt die Sprache
  basicstyle=\scriptsize\ttfamily,     % Setzt den Standardstil
  % keywordstyle=\color{red}\bfseries,    % Setzt den Stil für Schlüsselwörter
  identifierstyle=\color{blue},        % Identifier bekommen keine gesonderte formatierung
  commentstyle=\color{darkspringgreen},        % Stil für Kommentare
  stringstyle=\ttfamily,             % Stil für Strings (gekennzeichnet mit "String")
  breaklines=true,             % Zeilen werden umgebrochen
  numbers=left,                 % Zeilennummern links
  numberstyle=\tiny,             % Stil für die Seitennummern
  frame=single,                 % Rahmen
  % backgroundcolor=\color{myGrey},     % Hintergrundfarbe
  % caption={Java-Code},             % Caption
  tabsize=2                % Größe der Tabulatoren
}



\newcommand{\GETHeader}[6]{%
  \begin{titlepage}
    \pagestyle{empty} \enlargethispage*{25cm}\samepage{
      \vspace*{-2.5cm}
      \begin{center}
        %
        \begin{tabular}[b]{lr}
          %
          \hspace*{-1.5cm}
          \begin{tabular}{p{3.8cm}}
            \includegraphics[width=3.8cm]{./Figures/TUGraz}
          \end{tabular}
          %
          \hspace*{-1cm}
          %
          %
          \begin{tabular}{p{13.8cm}}
            \begin{flushright}
              \large
              Institut für Grundlagen und Theorie der Elektrotechnik\\
              ~\\
            \end{flushright}
          \end{tabular}
        \end{tabular}

        \vspace*{2.2cm}
        %
        \Huge {Elektrische Netzwerke und Mehrtore \\ Übung\\} %
        \vspace*{.5cm} \Large{ Wintersemester 2020\\}
        %
        \vspace*{1.5cm}

        \Huge{\textbf{#1}\\}

        %
        % \vspace*{0.8cm} \Large{Übungsdatum: {#2}\\}
        %
        \vspace*{1cm} \vfill
        %
        \Large{Gruppe: {#2}\\} \vspace*{0.5cm}%


        % \Large{Protokollführer(in): {#3}\\} \vspace*{1cm}

        \Large{Gruppenteilnehmer:\\} \vspace*{.1cm}
        % Name der beteiligten Studierenden
        \Large{#3} \vspace{1cm}
        %
        %
        \Large{Vortragende: #4\\} \vspace*{.1cm}   %\vspace*{1.5cm}
        % \Large{Betreuer(in): #6\\}
        \vspace*{1.5cm}
        %
        \Large{#5, am #6}
      \end{center}}%
    %
    \clearpage
  \end{titlepage}}








\begin{document}

\GETHeader                                                                              %  Bitte Ausfüllen!!!
% ----------------------------
{Protokoll Übung 3: \\ Schaltvorgang Kondensator}                         %  Übungstitel
% ----------------------------
% {25. Mai 2020}                                                                  %  Übungsdatum
% ----------------------------
{04}                                                                   %  Gruppen-Nr.
% ----------------------------
% {Matthias Fottner}                                                                      % Name des Protokollführers oder der Protokollführerin
% ----------------------------
{
  \begin{center}
    \begin{minipage}{0.28\linewidth}
      1. Matthias Fottner\\
      2. David Keller\\
      3. Moritz Woltron
    \end{minipage}
  \end{center}
}
% ----------------------------
{Helena Grabner}                                                                     %  Laborleiter(in)
% {Übung 2}                                                               %  Betreuer(in)
% ----------------------------
{Graz}                                                                                  %  Ort der Protokollerstellung
{\today}                                                                %  Datum Protokollerstellung


\newcommand{\unit}[1]{\,\text{#1}}


\tableofcontents

\newpage


\allowdisplaybreaks


\section{Ermitteln der DGL von $i_L(t)$ für $0 \leq t \leq 2\tau_1$} %David
\subsection{Schaltbild des Netzwerks für $0 \leq t \leq 2\tau_1$}%David

\begin{figure}[!htb]
	\begin{center}
		\begin{circuitikz} [european resistors, scale=1]
			
			
			\draw(0,0) to [L=$L$, v<={\color{blue}{$u_{L}$}}, i<={\color{red}{$i_L(t)$}}] (0,4)
			to [R=$R_1$, v<={\color{blue}{$u_{R1}$}}] ++ (4,0)
			to [vsource, v>={\color{blue}{$U_{0}$}}] ++ (0,-4)
			to [short] ++ (-4,0);
			%\draw[thin, <-, >=triangle 45] (2,2)node{$i_1$}  ++(-60:0.5) arc (-60:170:0.5);
			\draw [green,latex-] (2,2) node[above,color=green]{$m_1$} arc (-90:165:5mm);
			
			
			
			
			
			
			
%			\clip (-2,-5.7) rectangle (12.5, 5);
%			
%			\draw(0,0) to [R=$R_3$, v={\color{blue}{$U_{R3}$}}, i={\color{red}{$I_{R3}$}}, *-*] (4,0);
%			\draw(4,0) to [R=$R_4$, v={\color{blue}{$U_{R4}$}}, i={\color{red}{$I_{R4}$}}] (4,4);
%			\draw(4,0) to [R=$R_2$, v<={\color{blue}{$U_{R2}$}}, i<={\color{red}{$I_{R2}$}}] (4,-4);
%			
%			
%			
%			\draw(0,0) node[label={[font=\footnotesize]-190:n3}] {} to (0,4);
%			\draw(0,4) to[isource, i<={\color{red}{$I_{S2}$}}, -*] (4,4) node[label={[font=\footnotesize]90:n2}] {};
%			\draw(4,4) to [R=$R_5$, v={\color{blue}{$U_{R5} = 0\unit{A}$}}, i={\color{red}{$I_{R5} = 0\unit{A}$}}, -o] (8,4);
%			\draw(8,4) to (10,4);
%			
%			\draw(0,0) to[vsource, v<={\color{blue}{$U_{S1}$}}, i<={\color{red}{$I_{S1}^?$}}] (0,-4);
%			\draw(0,-4) to [R=$R_1$, v={\color{blue}{$U_{R1}$}}, i={\color{red}{$I_{R1}$}}, *-*] (4,-4) node[rground]{};
%			\draw(4,-4) to [short, -o] (8,-4);
%			\draw(8,-4) to (10,-4);
%			\draw(10,-4) to [C, v<={\color{blue}{$u_{C} = 0\unit{A}$}}, i<={\color{red}{$i_{C} = 0\unit{A}$}}] (10,4);
%			
%			\draw[FARROW, green!50!black] (-0.2, 0) to[bend left=-90, looseness=2.2] node[below left]{$U_{n3}$} (3.9,-4.2);
%			\draw[FARROW, green!50!black] (4.1, 3.95) to[bend right=-75, looseness=0.9] node[below right]{$U_{n4}$} (4.4,-3.9);
%			\draw[FARROW, green!50!black] (4.1, 0) to[bend right=-75, looseness=0.7] node[below right]{$U_{n1}$} (4.4,-3.7);
%			\draw[FARROW, green!50!black] (0.1, -3.8) to[bend right=-75, looseness=0.5] node[above]{$U_{n2}$} (3.9,-3.7);
			
			%\draw [brown] (current bounding box.south west) rectangle (current bounding box.north east);
			%\draw[gray,step=0.25] (-6,-3) grid (16.5, 11);
			
			
		\end{circuitikz}
	\end{center}
	\caption{Netzwerk mit allen eingezeichneten Strömen, (Knoten-)spannungen und Knoten}
	\label{fig:schaltplan_a}
\end{figure}

\subsection{Aufstellen der DGL mithilfe der allgemeinen Lösungsformel} %David
\section{Ermitteln der DGL von $i_L(t)$ für $t > 2\tau_1$}%Matthias
\subsection{Schaltbild des Netzwerks für $t > 2\tau_1$}%Matthias
\subsection{Aufstellen der Kirchhoff'schen Knoten- und Maschengleichungen}%Matthias
\subsection{Herleitung der DGL 2. Ordnung von $i_L(t)$ für $t > 2\tau_1$}%Matthias
\subsection{Interpretation der Parameter $\delta$,  $\omega_0$ und $\Omega_d$}%Matthias
\subsection{Anfangswertproblem}%Matthias
\subsubsection{Anfangsbedingungen}%Matthias
\subsubsection{Lösen von $K_1$ und $K_2$}%Matthias
\section{Plots und Simulationen} %David
\subsection{Matlab-Plot $i_L(t)$} %David
\subsection{PSpice-Plot $i_L(t)$ und $u_L(t)$} %David

\end{document}